% 七、问题三
\section{问题三：正式预报与滚动调整购电}
问题 3 在问题 2 之上扩展两件事：\textbf{信息}——每天 0:00、6:00、12:00、18:00 获得未来 24 小时整点光伏预报；\textbf{权限}——可在发布时刻调整尚未执行的购电量，调减部分付 50\% 违约费、调增部分按 1.5 倍电价购入。全天实际账单为四项分解
\begin{equation}
C_{\mathrm{real}}=\sum_t\Bigl[p_t\min(G^0_t,G^a_t)+0.5\,p_t\,(G^0_t-G^a_t)^+
+1.5\,p_t\,(G^a_t-G^0_t)^++5p_tE_t\Bigr],
\label{eq:q3bill}
\end{equation}
等价于 $\sum_tp_t[G^a_t+0.5|G^a_t-G^0_t|+5E_t]$：交付按最终执行量计价，偏离原计划的部分对称地付半价摩擦费。每个交付段只按最终提交量结算一次，被覆盖的中间临时决策不计费。

\subsection{预报接入与组合预测}
附件 3 发布后第 $h$ 小时的预报对应发布时刻加 $h$ 的整点值。以发布时刻已观测光伏为左锚点，把 24 小时整点预报插值到 10 分钟段末。0:00 的正式预报与历史预测（三日均值）各有所长，取凸组合
\begin{equation}
\widehat{V}=w\,\widehat{V}^{F}+(1-w)\,\widehat{V}^{H},\qquad w=0.388815,
\label{eq:q3combine}
\end{equation}
其中权重 $w$ 为最小化组合误差平方和的经验权重 $w=\sum e_H(e_H-e_F)/\sum(e_F-e_H)^2$（截断到 $[0,1]$），由一月内各历史发布时刻此前已完成的配对误差因果估计，2 月 1 日起冻结。6/12/18 时的预报误差显著更小，直接采用最新正式预报。滚动场景库按发布小时分组：每次只选\emph{同发布小时}、完整 24 小时已实现的最近 30 条三通道配对残差，跨午夜路径按精确结束时刻判断可用性。

\subsection{滚动重优化模型}
发布时刻 $\tau$（当天剩余 $q_\tau=144-6\tau$ 段）的重优化以当前真实储电量为初值，在冻结的 $G^0$ 上决策调整量。第一阶段变量为当天剩余段的调增 $\Delta^+_u$ 与调减 $\Delta^-_u$（跨场景共享），$G^a_u=G^0_u+\Delta^+_u-\Delta^-_u$；第二阶段为各场景储能追索、缺口补购与跨午夜次日暂拟购电 $\widetilde G_{\omega,u}$（按场景分开，只用于估计未来费用，不提交）。完整模型：
\begin{equation}
\begin{aligned}
\min_{\Delta^\pm,\{\cdot_\omega\}}\quad
& \Gamma_\tau+\sum_{u>\!6\tau}\bigl(1.5p_u\Delta^+_u-0.5p_u\Delta^-_u\bigr)
+\frac1M\sum_{\omega}\Bigl[\sum_{u\in\text{次日}}p_u\widetilde G_{\omega,u}
+\sum_{u}5p_uE_{\omega,u}\Bigr]-\frac1M\sum_\omega V(S_{\omega,\mathrm{end}})\\
\mathrm{s.t.}\quad
& G^a_u=G^0_u+\Delta^+_u-\Delta^-_u,\quad 0\le\Delta^-_u\le G^0_u,\ \Delta^+_u\ge0, && u>6\tau,\\
& G^a_u+E_{\omega,u}+V_{\omega,u}+D_{\omega,u}=L_{\omega,u}+C_{\omega,u}+W_{\omega,u}, && \forall\omega,\ u>6\tau,\\
& \widetilde G_{\omega,u}+E_{\omega,u}+V_{\omega,u}+D_{\omega,u}=L_{\omega,u}+C_{\omega,u}+W_{\omega,u}, && \forall\omega,\ u\in\text{次日},\\
& S_{\omega,u}=S_{\omega,u-1}+\eta C_{\omega,u}-D_{\omega,u}/\eta,\quad S_{\omega,6\tau}=S_{\mathrm{now}}, && \forall\omega,u,\\
& 1200\le S_{\omega,u}\le10800,\ 0\le C_{\omega,u},D_{\omega,u}\le\bar P,\ E,W,\widetilde G\ge0, && \forall\omega,u,
\end{aligned}
\label{eq:q3lp}
\end{equation}
其中 $\Gamma_\tau$ 为按式\eqref{eq:q3bill} 结算的当天已执行部分费用（常数），$V(\cdot)$ 为次日价值函数（主方法；24h 基线用 $-vS$）。调增与调减在同段的净费用为 $+p_u>0$，故最优解中二者不共存。前瞻始终为完整 24 小时（如 18:00 覆盖当天 18--24 时与次日 0--18 时），但\textbf{只提交到下一发布时刻}：之后的段是临时决策，下轮重优化以更新的信息覆盖之。执行层与问题 2 完全相同——每轮重优化后按新的 $\bar H$ 与保留水平因果执行到下一发布时刻。

\begin{proposition}[重优化的代理不劣性]\label{prop:readjust}
固定发布时刻 $\tau$ 的信息状态（同一场景集、同一未来预测、同一当前储电量与冻结计划 $G^0$）。式\eqref{eq:q3lp} 的重优化线性规划的最优值不超过不调整方案（$\Delta^\pm\equiv0$，即 $G^a=G^0$）的代理费用；差值即该发布时刻在代理口径下的信息价值，非负。
\end{proposition}

\begin{proof}
取 $\Delta^+=\Delta^-=0$，其余变量取不调整方案的第二阶段最优追索值，则该组取值满足式\eqref{eq:q3lp} 的全部约束，是其可行解，目标值恰为不调整方案的代理费用。线性规划最优值不超过任一可行解的目标值，命题得证。
\end{proof}


命题~\ref{prop:readjust} 给出的是\emph{代理口径}的方向保证：更多发布时刻不会让重优化 LP 变差。实际账单由真值结算，不被该命题覆盖，需由全年实验回答（\S\ref{sec:q3issues}）。数值上，四个指定日 12:00 重优化的代理改善最小值为 56.70 元 $\ge0$，与命题一致。

图~\ref{fig:fig_f08_rolling_timeline} 给出滚动机制的时间结构：每轮 24 小时前瞻、只提交一个发布间隔、临时决策可被覆盖。

\begin{figure}[!htbp]\centering
\safeincludegraphics[width=0.94\textwidth]{fig_f08_rolling_timeline}
\caption{滚动重优化时间线：四个发布时刻各前瞻完整 24 小时（可跨午夜），实际只提交到下一发布时刻（实色段），其后为可被覆盖的临时决策（虚框段）}
\label{fig:fig_f08_rolling_timeline}
\end{figure}

\subsection{是否需要引入其他时刻的预报}\label{sec:q3issues}
固定预测方法、场景构造、终端定价与结算口径，逐级开放发布时刻集合 $\{0\}\subset\{0,6\}\subset\{0,6,12\}\subset\{0,6,12,18\}$，四组配置均连续执行全年：

\begin{table}[!htbp]
\centering\small
\caption{问3发布时刻组合与相邻组合的费用变化}\label{tab:issues}
\begin{tabular}{lrrr}
\toprule
发布时刻/h & 总费用/元 & 紧急费/元 & 相邻节省/元 \\
\midrule
0 & 13351182.68 & 579245.64 & --- \\
0,6 & 13312853.41 & 353878.30 & 38329.27 \\
0,6,12 & 13075459.34 & 265785.37 & 237394.07 \\
0,6,12,18 & 13024596.85 & 192075.49 & 50862.49 \\
\bottomrule
\end{tabular}
\end{table}


相邻差即该发布时刻的边际价值（在此固定增加顺序下）：新增 6:00 全年节省 3.83 万元，新增 \textbf{12:00 节省 23.74 万元}，新增 18:00 节省 5.10 万元（图~\ref{fig:fig_f09_issue_marginal}）。12:00 的边际贡献为何一枝独秀？两个因素在此交汇：其一，光伏不确定性集中在午间出力峰值段，12:00 预报恰在峰值前夕把当天剩余光伏路径的方差大幅压缩；其二，12:00 时当天还有一半的段未执行，调整窗口足够长，信息可以兑换成实际的采购修正。相比之下，6:00 预报虽早，但此时距光伏峰值尚远、预报本身的增量信息少；18:00 预报虽准，但当天只剩 6 小时可调，信息变现空间小。\textbf{结论：需要引入其他时刻的预报，其中 12:00 预报的引入最为关键}；若受成本限制只能保留一个日内更新时刻，应保留 12:00。

\begin{figure}[!htbp]\centering
\safeincludegraphics[width=0.86\textwidth]{fig_f09_issue_marginal}
\caption{逐级增加发布时刻的全年边际节省：12:00 预报的边际价值（23.74 万元/年）约为 6:00 与 18:00 之和的三倍——信息增量与剩余调整窗口在午间交汇}
\label{fig:fig_f09_issue_marginal}
\end{figure}

\subsection{结果}
主方法（48h+$V(e)$）全年实际费用 \textbf{1298.96 万元}，比问题 2 主方法降低 54.01 万元（4.0\%）：这是正式预报信息与日内调整权限的联合价值。四项费用分解中紧急购电费降至 19.16 万元（问题 2 为 78.23 万元）——日内调整把大部分本须紧急补购的缺口提前转化为 1.5 倍价的调增，避开了 5 倍价；紧急购电量占比从 0.70\% 降至 0.17\%。24 小时基线为 1302.46 万元，与问题 2 一致地支持 48h 组件。月度与指定日结果见 \S\ref{sec:results}。
